% =====================================================================
%  Preface -- description of the course
% =====================================================================
\chapter*{Preface}
\addcontentsline{toc}{chapter}{Preface}
\markboth{Preface}{Preface}

This book accompanies the course \emph{Evaluation, Deployment and Monitoring of Models}
(EDM), taught in the third year of the Degree in Data Science (\emph{Grado en Ciencia de
Datos}, GCD) at the Universitat Polit\`ecnica de Val\`encia. By the time students reach
this course they already know how to build predictive models. This course is about
everything that happens around the model: how to know whether it is good, how to combine
it with other models and with prior knowledge, how to make it reliable, how to put it into
production, and how to keep it working once the world starts changing.

\section*{What the course is about}

The course addresses the evaluation, deployment and maintenance of predictive models and
recommender systems, exploiting jointly the knowledge induced from data and the prior
knowledge of the domain (constraints, costs, manual models). Its main topics are:
\begin{itemize}
  \item evaluation procedures and metrics, and the calibration of probabilistic models;
  \item automated machine learning (AutoML) and the fusion and combination of models;
  \item reliable models: fairness, causality, explainability and security;
  \item the deployment of models and their monitoring and review in production.
\end{itemize}

\section*{Organisation}

The book follows the seven units of the syllabus (\cref{tab:units}). Each unit is a chapter
that can be read on its own, although later units frequently rely on the evaluation tools
introduced in \cref{ch:evaluation}.

\begin{table}[htbp]
\centering\small
\caption{Units of the course.}\label{tab:units}
\begin{tabularx}{\textwidth}{@{}l>{\raggedright\arraybackslash}p{3.6cm}X@{}}
\toprule
& Unit & Main contents\\
\midrule
1 & Model evaluation & Evaluation protocols; metrics for crisp, scoring and probabilistic
classifiers; ROC analysis; calibration; imbalanced problems; cost-based evaluation;
regression; specialised tasks; unsupervised learning; statistical comparison.\\
2 & Hybrid models & Wisdom of crowds; combining models; ensembles; hybrid approaches;
cautious models and delegating classifiers; performance evaluation and drawbacks.\\
3 & Reliable models & Fairness in machine learning; causality; interpretable models,
model-agnostic and example-based explanations; machine learning and security, adversarial
learning and robustness.\\
4 & Learning to learn & Automating data science (AutoML); inductive and genetic
programming; process mining and robotic process automation (RPA).\\
5 & Deployment of models & Choosing a deployment option; development infrastructure;
ModelOps; building the technology stack.\\
6 & Monitoring and updating models & What to monitor; CI/CD for machine learning; A/B testing
and experimentation (groups, duration, estimating the better variant); infrastructure.\\
7 & Applications & Recommender systems (content-based, collaborative and hybrid); customer
relationship management (campaign design, mailings, rankings).\\
\bottomrule
\end{tabularx}
\end{table}

The practical part of the course consists of ten laboratory sessions: model evaluation
(three sessions), hybrid models (one), reliable models (three), learning to learn (one),
deployment of models (one) and applications (one).

\section*{How to use this book}

Definitions, worked examples and code are interleaved with the explanations. Several kinds of
boxes highlight special content:
\begin{description}[font=\sffamily\bfseries\small]
  \item[Key idea] a principle worth remembering after the details are forgotten;
  \item[Example] a worked example, usually with all the numbers so that it can be checked
        by hand;
  \item[Pitfall] a frequent mistake, often one that produces optimistic estimates;
  \item[Going further] material that goes beyond the core of the course.
\end{description}

All the code is written in Python using \textsf{scikit-learn} \citep{pedregosa2011scikit},
\textsf{NumPy}, \textsf{pandas}, \textsf{SciPy}, \textsf{matplotlib} and
\textsf{imbalanced-learn} \citep{lemaitre2017imbalanced}. Every listing has been executed and
the output printed below it is the real output (it was produced with Python~3.12, \textsf{scikit-learn}~1.8.0,
\textsf{NumPy}~2.4.4, \textsf{SciPy}~1.17.1, \textsf{pandas}~3.0.2 and \textsf{imbalanced-learn}~0.14.2;
other versions or platforms may differ in the last decimals). The
companion folder \code{code/} contains the listings and the script that generates every figure
of the book; all random seeds are fixed, so the figures and the numbers quoted in the text
can be reproduced in that environment, up to small numerical differences across platforms. Each unit ends with a summary, bibliographic notes and exercises,
with solutions to most of them.

\section*{Recommended reading}

The following books cover, between them, the whole course. For the foundations of
statistical learning, \citet{hastie2009elements} is the standard reference, and
\citet{flach2012machine} offers a modern view of machine learning with special emphasis on
evaluation and ROC analysis; \citet{sammut2010encyclopedia} provides short, authoritative
entries on most topics. \citet{japkowicz2011evaluating} is entirely devoted to the evaluation
of classifiers. \citet{provost2013data} connects data science with business decisions and
\citet{segaran2007programming} introduces collective intelligence and recommender systems.
Reliable models are covered by \citet{molnar2022interpretable} (interpretability),
\citet{barocas2023fairness} (fairness) and \citet{chio2018machine} (security).
\citet{kyriakides2019hands} is a practical introduction to ensembles, and
\citet{ameisen2020building}, \citet{kelleher2019machine} and \citet{burkov2020machine} deal
with building, deploying and maintaining machine learning systems.

\bigskip
\noindent\emph{This is a draft of the complete book: all seven units, with their code, figures, exercises and
solutions. Comments and corrections are welcome.}
